%%%%%%%%%%%%%%%%%%%%%%%%%%%%%%%%%%%%%%%%%%%%%%%%%%%%%%%%%%%%%%%%%%%%%%%%%%%%%%
%  mm-howtouse.tex
%  Editorial conventions for the series.  Shared, identical, across all five
%  volumes so that a reader who learns the apparatus once can use any volume.
%%%%%%%%%%%%%%%%%%%%%%%%%%%%%%%%%%%%%%%%%%%%%%%%%%%%%%%%%%%%%%%%%%%%%%%%%%%%%%

\frontchapter{How to Use This Study Text}

\section*{The shape of a volume}
\addcontentsline{toc}{section}{The shape of a volume}

Each of the five volumes in this series covers exactly one unit of the
IM345AT syllabus and is self-contained. Every volume opens with the official
syllabus for the whole course, so that the unit in hand can always be located
within it, and closes with a set of appendices carrying the concordance,
the question bank, the glossary and a quick-reference sheet.

The numbered chapters of each volume map one-to-one onto the topic phrases of
that unit's official syllabus descriptor. The order of the chapters is the
order of the syllabus. Appendix~A sets out that mapping explicitly, phrase by
phrase, so coverage can be audited rather than assumed.

\section*{The five panel types}
\addcontentsline{toc}{section}{The five panel types}

Material that is not continuous prose is set in one of five panels. Each has a
fixed meaning throughout the series.

\begin{definitionbox}[the term being defined]
A formal definition, worth reproducing close to verbatim in an answer script.
The term appears in the panel heading and again in the index.
\end{definitionbox}

\begin{keybox}[High Yield]
A result, list or distinction that carries disproportionate examination weight.
Where a topic has recurred across papers, the recurrence is recorded as a
bracketed tag. \pyq{CIE-I, 3 Apr 2025; CIE-I, 16 Apr 2026}
\end{keybox}

\begin{exambox}[Model Answer]
An answer written at the length and in the register the marking scheme
rewards. Where an official scheme of valuation was released, the panel follows
it.
\end{exambox}

\begin{examplebox}[Worked Example]
A concrete illustration --- a named company, a named persona, a worked
calculation --- rather than a general statement.
\end{examplebox}

\begin{warnbox}
A common and costly error: a distinction usually collapsed, a term usually
misused, or an answer that looks complete and is not.
\end{warnbox}

A sixth panel, in green, sets out a framework or process as a single
sequence:

\begin{frameworkbox}[Framework]
Step one \ra\ step two \ra\ step three \ra\ back to step one.
\end{frameworkbox}

\section*{How previous-year questions are recorded}
\addcontentsline{toc}{section}{How previous-year questions are recorded}

\begin{keybox}[The Bracket Convention]
A past-paper appearance is recorded \textbf{only} as a bracketed source tag
attached to the material, never as a sentence of narrative. So the text reads
\begin{quote}\sffamily\footnotesize
Push versus pull marketing. \pyq{CIE-I, 3 Apr 2025; CIE-I, 16 Apr 2026; SEE, Jun/Jul 2025 --- 2 M}
\end{quote}
and never ``this was asked in 2025 and again in 2026''. The tag carries the
paper, the date and, where the mark allocation is known, the marks. Multiple
appearances are separated by semicolons.
\end{keybox}

\noindent The abbreviations used inside tags are:

\vspace{4pt}
{\sffamily\small
\begin{tabular}{@{}P{0.14\textwidth}P{0.80\textwidth}@{}}
\toprule[1.1pt]
\rlab{CIE-I} & First Continuous Internal Evaluation test\\
\rlab{CIE-II} & Second Continuous Internal Evaluation test\\
\rlab{CIE-III} & Third Continuous Internal Evaluation test\\
\rlab{SEE} & Semester End Examination\\
\rlab{2 M} & Mark allocation of the question as printed on the paper\\
\bottomrule[1.1pt]
\end{tabular}}

\vspace{10pt}

Appendix~B of each volume reproduces every past-paper question traced to that
unit, in full, grouped by mark weight, each carrying the same bracketed tag.

\section*{Cross-references, figures and the index}
\addcontentsline{toc}{section}{Cross-references, figures and the index}

Cross-references are given as \S\,chapter.section --- so \S\,1.4 is the fourth
section of Chapter~1. Every figure in this series is drawn in \TeX{} rather
than imported as an image, so the diagrams print at full resolution at any
size and use the same typeface as the surrounding text. Figures and tables are
listed after the table of contents.

Terms set in \textbf{\textit{bold italic}} on first use are indexed. The index
at the back of each volume is a real back-of-book index built from those
terms, not a repeat of the contents list.

\section*{Status of this text}
\addcontentsline{toc}{section}{Status of this text}

\begin{warnbox}[Unofficial Document]
This is a student-compiled study text. It is not issued, endorsed or verified
by the Department of Industrial Engineering and Management or by RV College of
Engineering. The syllabus reproduced in the front matter, the assessment
rubrics and the previous-year questions are transcribed from official
documents; the explanatory prose around them is not official. Where the source
lecture notes were incomplete, gaps have been filled from the prescribed
reference texts listed in the front matter. Always cross-check against the
official syllabus and your own class notes.
\end{warnbox}
